% Appendix D, typeset from project/register_map.md and project/spi_register_protocol.md.

\chapter{The register map and the SPI protocol}
\label{app:protocol}

This is the contract at the middle of the project, and it is written for \textbf{two classes at
once}: this course builds the FPGA side as an SPI slave (\chapref{c18:ch} and \ref{c19:ch}), and the
parallel class builds the microcontroller side as master. Both halves are written against this
document, not against each other's code. \textbf{If an implementation and this appendix disagree, the
implementation is wrong.}

The contract says nothing about CAN as a protocol; it assumes it. \canbook{} is the description, and
its chapter~7 is the one explaining why this layer is needed at all: a controller signalling with
one-cycle pulses cannot be polled over a serial bus.

Two things are kept apart, and it is worth knowing which is which. The \emph{register map}
(\secref{app:protocol:map} to \ref{app:protocol:order}) says \textbf{what} the registers mean. The
\emph{transaction format} (\secref{app:protocol:electrical} to \ref{app:protocol:timing}) says
\textbf{how} they are reached over the wires.

\begin{warning}
An older version of this project had the driver reach the registers via memory-mapped I/O at the base
address \code{0xFF200000}. That assumes a processor with an address bus out to the FPGA fabric, which
the microcontroller in this system does not have. The registers are reached over SPI, and the
\textbf{Index} column below, not the offset, is what goes in the command byte.
\end{warning}

% ------------------------------------------------------------------------------------------
\section{The registers}
\label{app:protocol:map}

\begin{booktable}{clllL}
  \thead{Index} & \thead{Register} & \thead{Offset} & \thead{Access} & \thead{Description} \\
  \midrule
  0 & \code{STATUS} & \code{0x00} & R & Bit 0: TX ready. Bit 1: RX valid. Bit 2: Error. \\
  1 & \code{TX_ID} & \code{0x04} & R/W & 11-bit transmit ID (bits 10--0). \\
  2 & \code{TX_DLC} & \code{0x08} & R/W & Data Length Code (bits 3--0, value 0--8). \\
  3 & \code{TX_DATA_LO} & \code{0x0C} & R/W & Transmit data bytes 0--3 (byte 0 = MSB). \\
  4 & \code{TX_DATA_HI} & \code{0x10} & R/W & Transmit data bytes 4--7 (byte 4 = MSB). \\
  5 & \code{TX_SEND} & \code{0x14} & W & Write \code{0x1} to trigger a transmission. \\
  6 & \code{RX_ID} & \code{0x18} & R & Received identifier (bits 10--0). \\
  7 & \code{RX_DLC} & \code{0x1C} & R & Received DLC (bits 3--0). \\
  8 & \code{RX_DATA_LO} & \code{0x20} & R & Received data bytes 0--3. \\
  9 & \code{RX_DATA_HI} & \code{0x24} & R & Received data bytes 4--7. \\
  10 & \code{RX_ACK} & \code{0x28} & W & Write \code{0x1} to acknowledge and empty the RX
    buffer. \\
  11 & \code{ERROR_FLAGS} & \code{0x2C} & R/W & The error register; write \code{0x0} to clear. \\
  12 & \code{TX_ABORT} & \code{0x30} & W & Write \code{0x1} to force TX ready back after an aborted
    transmission. \\
\end{booktable}

Indices 13--15 are reserved: a read gives \code{0x00000000}, a write is ignored.

\code{Index} is \code{Offset / 4}. The offset is still in the table because the driver's constants are
written in that form, but it is the index that goes out on the wire.

% ------------------------------------------------------------------------------------------
\section{\code{STATUS}}
\label{app:protocol:status}

\begin{booktable}{cLLL}
  \thead{Bit} & \thead{Meaning} & \thead{Set by} & \thead{Cleared by} \\
  \midrule
  0 & TX ready & reset; the controller's \code{tx_done} pulse; a rising edge on \code{error} while
    the bit is low; an accepted write to \code{TX_ABORT} & an accepted write to \code{TX_SEND} \\
  1 & RX valid & the controller's \code{rx_valid} pulse, which also latches \code{RX_*} & a write of
    \code{0x1} to \code{RX_ACK} \\
  2 & Error & a \textbf{rising edge} on the controller's \code{error} level & a write of \code{0x0}
    to \code{ERROR_FLAGS} \\
\end{booktable}

The point of the whole layer is in that table. The controller signals with one-cycle pulses, 20 ns at
50 MHz. A driver polling over SPI sees the system a few microseconds at a time at best, so a pulse is
gone thousands of cycles before the next poll arrives. The register bank turns the pulses into
\textbf{sticky, pollable levels with an explicit clear}.

% ------------------------------------------------------------------------------------------
\section{Semantics, register by register}
\label{app:protocol:semantics}

\begin{itemize}
  \item \textbf{\code{TX_SEND}} triggers a single \code{tx_req} pulse when a value with \textbf{bit 0
    set} is written and \code{STATUS} bit 0 is set. The other bits of the written word are ignored.
    If \code{STATUS} bit 0 is low the write is ignored: a transmission is in progress. A write with
    bit 0 low does nothing. A read gives \code{0x00000000}, since the register stores nothing but
    \emph{does} something. Writes are edge events at this layer, so a repeated write is a repeated
    request, never a held level.
  \item \textbf{\code{RX_ACK}} clears \code{STATUS} bit 1. Requires bit 0 set; other bits are
    ignored. A read gives \code{0x00000000}.
  \item \textbf{\code{TX_ABORT}} sets \code{STATUS} bit 0 back to \code{1}, and does nothing else.
    Only a write with \textbf{bit 0 set} does anything; a read gives \code{0x00000000}. The register
    does not reach the controller \textemdash{} a controller that has aborted a transmission is
    already back at idle and can transmit again; it is only the bank's sticky bit that is stuck. This
    is the escape hatch: in a correct design it should never be needed, since bit 0 already has two
    automatic routes back. A driver that finds it needs this has found a bug worth reporting.
  \item \textbf{\code{ERROR_FLAGS}} is a latch-and-clear-on-zero construct, not a general register:
    only a write where \textbf{the whole word is zero} clears. \code{0x000000FF} therefore does not
    clear, even though bit 0 is set. Bit 0 is the latch and mirrors \code{STATUS} bit 2; bits 31--1
    read as \code{0}. The name is plural for historical reasons.
  \item \textbf{Masking on write.} \code{TX_ID} stores bits 10--0 and \code{TX_DLC} bits 3--0;
    surplus bits read back as zero. A write of \code{0xFFFFFFFF} therefore reads back as
    \code{0x000007FF} and \code{0x0000000F} respectively. \code{TX_DATA_LO}/\code{HI} store all 32
    bits.
  \item \textbf{No validation.} A \code{TX_DLC} of \code{0xF} is stored and passed on to the
    controller as it is. Rejecting \code{dlc > 8} is the driver's responsibility. Doing it on both
    sides is exactly how the two copies of the rule drift apart.
  \item \textbf{No buffering on the receive side.} If a new frame arrives before the previous has
    been acknowledged it is overwritten, and the newest wins. That is a deliberate, documented choice
    and not a bug; a real controller has a receive queue.
\end{itemize}

\subsection{Details a driver needs}
\label{app:protocol:driver}

The points below belong to the contract as much as the tables above.

\begin{itemize}
  \item \textbf{Every register is 32 bits wide.} All five bytes of a transaction are always
    transferred, even for a register with only four meaningful bits.
  \item \textbf{\code{STATUS} bits 31--3 read as \code{0}.} The bank's read multiplexer zeroes
    everything above bit 2, so a comparison against \code{0x1} is safe.
  \item \textbf{\code{ERROR_FLAGS} does not distinguish causes.} The four causes \textemdash{} lost
    arbitration, stuffing error, failed CRC and a received remote frame \textemdash{} all give the
    same single bit.
  \item \textbf{\code{RX_DATA_LO}/\code{HI} beyond \code{RX_DLC} read as \code{0x00}.} The controller
    clears its data accumulator at every frame start, so bytes above the received length never carry
    residue from an earlier frame: a frame with \code{DLC = 2} following one with \code{DLC = 8} reads
    as two bytes of data and six zeros. \textbf{Mask against \code{RX_DLC} all the same}
    \textemdash{} the zeros are the controller's guarantee, not the frame's contents, and a driver
    copying eight bytes reports six bytes that were never sent.
  \item \textbf{\code{TX_DATA_LO}/\code{HI} beyond \code{TX_DLC} are never sent}, so they need not be
    cleared.
  \item \textbf{Writes to \code{STATUS}, \code{RX_ID}, \code{RX_DLC} and
    \code{RX_DATA_LO}/\code{HI} are silently ignored.} The transaction completes normally; no error is
    reported.
  \item \textbf{There is no timeout to derive.} The map does not say how long a transmission is, so a
    driver that wants to give up has to choose its own timeout. A guide value: a maximum frame is
    about 130 bits, and at 1 Mbit/s therefore about 130 µs, plus waiting for an idle bus.
  \item \textbf{A lost arbitration restores \code{STATUS} bit 0}, exactly as a successful
    transmission does. A driver therefore does not have to distinguish the outcomes in order to
    transmit again: poll bit 0, and read \code{ERROR_FLAGS} to find out how it went.
\end{itemize}

\begin{aside}
\textbf{What ``the transmission is over'' means here.} \code{STATUS} bit 0 coming back says that the
transmit port is free again, not that the frame was delivered. A lost arbitration gives back the same
bit as a clean transmission. \code{ERROR_FLAGS} is what tells them apart, and the ordinary polling
loop is therefore: wait for bit 0, then read \code{ERROR_FLAGS} to decide whether the frame should be
resent. \Secref{app:project:limits} lists the other simplifications against real CAN.
\end{aside}

% ------------------------------------------------------------------------------------------
\section{Data byte order}
\label{app:protocol:order}

\code{tx_data} and \code{rx_data} are 64 bits wide, with data byte 0 in the most significant bits:

\begin{widedrawing}
bit 63                                                              bit 0
+--------+--------+--------+--------+--------+--------+--------+--------+
| byte 0 | byte 1 | byte 2 | byte 3 | byte 4 | byte 5 | byte 6 | byte 7 |
+--------+--------+--------+--------+--------+--------+--------+--------+
 \___________ TX_DATA_LO ___________/ \___________ TX_DATA_HI __________/
\end{widedrawing}

The names \code{LO} and \code{HI} refer to the register pair's order, not to bit significance.
\code{TX_DATA_LO} carries the \emph{first} four data bytes.

% ------------------------------------------------------------------------------------------
\section{Counterparts in the hardware}
\label{app:protocol:hw}

\begin{booktable}{LL}
  \thead{Register} & \thead{Port on \code{can_controller}} \\
  \midrule
  \code{TX_ID}, \code{TX_DLC}, \code{TX_DATA_LO}/\code{HI} & \code{tx_id}, \code{tx_dlc},
    \code{tx_data} \\
  \code{TX_SEND} & \code{tx_req} (a pulse) \\
  \code{STATUS} bit 0 & \code{tx_done} (a pulse), latched into a level \\
  \code{RX_ID}, \code{RX_DLC}, \code{RX_DATA_LO}/\code{HI} & \code{rx_id}, \code{rx_dlc},
    \code{rx_data} \\
  \code{STATUS} bit 1 & \code{rx_valid} (a pulse), latched into a level \\
  \code{STATUS} bit 2, \code{ERROR_FLAGS} & \code{error} (a level), latched on a rising edge \\
\end{booktable}

The translation between the two columns is exactly what \code{register_bank} does, and nothing more.

% ------------------------------------------------------------------------------------------
\section{Electrical parameters and framing}
\label{app:protocol:electrical}

\begin{booktable}{lL}
  \thead{Parameter} & \thead{Value} \\
  \midrule
  Roles & The microcontroller is SPI master, the FPGA slave. \\
  Mode & SPI mode 0 (CPOL = 0, CPHA = 0): SCK idles low, both sides sample on the rising edge. \\
  Bit order & MSB first, in every byte. \\
  SCK frequency & $\le$ 1 MHz (the AVR32DB28 runs at 4 MHz and divides by 4). See
    \secref{app:protocol:timing} for why. \\
  \code{SS} & Active low. Framing signal: one transaction per low period. \\
  Voltage & The AVR32DB28's core runs at 5 V, the DE0-CV at 3.3 V, and no wire passes through a level
    shifter. The four SPI wires sit on \code{PORTC}, the AVR DB's second supply domain, with
    \code{VDDIO2} supplied from the DE0-CV's 3.3 V rail. \\
\end{booktable}

\subsection{Pin assignment}
\label{app:protocol:pins}

Both sides are fixed here. The microcontroller side uses \code{SPI0}'s \textbf{ALT1} pin mux on the
AVR32DB28, which puts the peripheral on \code{PC0}--\code{PC3}; the FPGA side is free, and this
section spends that freedom on making the boards interchangeable.

ALT1 is no arbitrary choice. \code{PORTC} is the AVR DB's \textbf{MVIO} domain: it is supplied from
\code{VDDIO2} rather than from \code{VDD}, so tying \code{VDDIO2} to the DE0-CV's 3.3 V rail puts all
four SPI wires in a 3.3 V domain while the core goes on running at 5 V. Every level the FPGA sees is
then a level it is specified for, and the 3.3 V the FPGA drives back is read against a 3.3 V threshold
rather than a 5 V one. On \code{SPI0}'s default mux the same four signals would sit on
\code{PA4}--\code{PA7} in the \code{VDD} domain, and the bench would need a level shifter to be safe.
That is why the mux is part of the contract rather than every group's own choice.

The microcontroller side therefore owes two things beyond the pin numbers, and both belong to the
driver class: \code{PORTMUX.SPIROUTEA} has to select \code{ALT1}, and \code{VDDIO2} has to be
powered before \code{PORTC} does anything at all. \code{MVIO.STATUS} reports whether it is.

\begin{warning}
\textbf{This constrains the board, not just the firmware.} The hat carrying the AVR32DB28 is built
for this system, and the pin table below places three requirements on its layout.
\code{PC0}--\code{PC3} carry SPI0 and nothing else \textemdash{} no LED, no button, and above all no
debouncing capacitor on \code{SS}, which at 1 MHz would make the framing edge unusable. Nothing else
sits anywhere on \code{PORTC}, since the whole port is a 3.3 V domain as soon as \code{VDDIO2} is.
And \code{VDDIO2} is routed as a net of its own out to a pin, decoupled with 100 nF near the package,
rather than tied to \code{VDD} on the board. None of the three can be fixed in firmware once the board
is etched.

\textbf{A single fuse decides whether any of the above is true.} MVIO's dual-supply mode is selected
by \code{MVSYSCFG} in \code{FUSE.SYSCFG1}. In single-supply mode \code{VDDIO2} is tied internally to
\code{VDD}, \code{PORTC} becomes a 5 V port, and the wiring in this section then drives 5 V into FPGA
inputs that cannot take 5 V \textemdash{} while the bench looks exactly like a correct one. Confirm
the fuse against the datasheet before the first node is connected, not after the first board has
behaved oddly.
\end{warning}

\begin{booktable}{llll}
  \thead{Signal} & \thead{AVR32DB28} & \thead{DE0-CV port} & \thead{DE0-CV pin} \\
  \midrule
  \code{SCK} & \code{PC2} & \code{sclk} & \code{GPIO_0(4)} \\
  \code{MOSI} & \code{PC0} & \code{mosi} & \code{GPIO_0(5)} \\
  \code{MISO} & \code{PC1} & \code{miso} & \code{GPIO_0(6)} \\
  \code{SS} & \code{PC3} & \code{ss} & \code{GPIO_0(7)} \\
  I/O supply & \code{VDDIO2} & \textemdash{} & 3.3 V \\
  Ground & \code{GND} & \textemdash{} & \code{GND} \\
\end{booktable}

\code{GPIO_0(0)} to \code{GPIO_0(2)} are reserved for the CAN side \textemdash{} \code{tx_bus},
\code{bus_en} and the wired-AND bus line \textemdash{} and \code{GPIO_0(3)} is deliberately left
unused as a gap between the two groups, so that a cable loom sitting one pin off does not short
anything between the domains.

\textbf{Why fix them at all, when Quartus does not care?} Nothing in the design breaks if a group
picks other pins; the point is what happens at the bench. With an assignment shared by both classes,
any microcontroller node fits any FPGA board: a cable loom can be built once instead of per pair, and
a software group whose partner board is not ready can borrow a working one instead of falling back on
simulation. That last case is not hypothetical \textemdash{} the two courses hold their bring-up and
their final presentations the same week.

A group \emph{may} use other pins, but only if the board and the node it is paired with agree, and
only if it is written down where the other class can see it. An undocumented pin change shows up as a
node that answers register reads and never sends a frame, which is an expensive way to discover a
wiring convention.

% ------------------------------------------------------------------------------------------
\section{Transactions}
\label{app:protocol:transactions}

Every transaction is exactly \textbf{5 bytes} long: one command byte, then four data bytes. \code{SS}
falls before the command byte and rises after the fifth byte; it has to stay low for the whole
transaction.

\begin{textdrawing}
Bit:      7    6    5    4    3    2    1    0
        +----+----+----+----+----+----+----+----+
        | W  | 0  | 0  | 0  |   register index  |
        +----+----+----+----+----+----+----+----+
\end{textdrawing}

\begin{itemize}
  \item \textbf{Bit 7 (\code{W})}: \code{1} = write, \code{0} = read.
  \item \textbf{Bits 6--4}: reserved, must be \code{0}.
  \item \textbf{Bits 3--0}: the register index, \code{offset / 4} from \secref{app:protocol:map}
    (0--12).
\end{itemize}

\paragraph{Write (\code{W = 1}).}
The master sends the register value in the four data bytes, \textbf{MSB first} (bits 31--24 in the
first data byte). The slave commits the assembled value to the addressed register only once the fifth
byte is complete. Whatever the slave drives on \code{MISO} during a write is meaningless; the master
discards it.

\paragraph{Read (\code{W = 0}).}
At the end of the command byte the slave \textbf{latches} the addressed register's current value
\textbf{once}, into its response shift register. The four data bytes then clock that value out on
\code{MISO}, MSB first, while the master sends dummy bytes \code{0x00}. The latch-once rule is not an
implementation detail but part of the contract: the four bytes always belong to \emph{one} coherent
sample of the register, taken at a single instant \textemdash{} \code{STATUS} in particular must never
be stitched together from two separate moments.

\paragraph{Abort.}
If \code{SS} rises before the fifth byte is complete, the transaction is abandoned \textbf{with no
side effects}: nothing is committed, no trigger fires, and the slave returns to idle, ready for the
next falling edge on \code{SS}. That is what makes the slave self-healing after a glitch or a master
that resets.

\paragraph{Invalid commands.}
A read of a reserved index (13--15) gives \code{0x00000000}; a write to one is ignored. A command byte
with any reserved bit (6--4) set is invalid in the same way: no read is latched and no write is
committed, and the slave still consumes the whole five-byte frame before returning to idle. Neither
case is an error the slave reports \textemdash{} there is no error channel at this layer. Enforcing
the rule about reserved \emph{command} bits is \code{spi_reg_bridge}'s job; \code{register_bank} never
sees more than a 4-bit index.

% ------------------------------------------------------------------------------------------
\section{Timing requirements}
\label{app:protocol:timing}

\begin{itemize}
  \item The slave \textbf{never} uses \code{SCK} as a clock. All three inputs (\code{SCK},
    \code{MOSI}, \code{SS}) are synchronised into the 50 MHz domain with two flip-flops each, and
    \code{SCK} edges are \emph{detected}, not clocked on (\chapref{c4:ch} for why, \chapref{c18:ch}
    for the module that does it). \code{spi_slave} builds its synchroniser in place rather than
    instantiating \code{meta_prev}, so that \code{spi_slave} depends on nothing in
    \file{controller/}.
  \item Two-flip-flop synchronisation plus edge detection costs a few 50 MHz cycles per edge; at
    SCK $\le$ 1 MHz there are $\ge$ 50 system cycles per SPI bit, an order of magnitude of margin. The
    protocol therefore guarantees correct operation only up to 1 MHz; faster may work, and is
    deliberately not promised.
  \item \code{MISO} changes on \code{SCK}'s falling edge and is sampled by the master on the rising
    one, as mode 0 prescribes.
  \item The master has to raise \code{SS} between transactions. Back-to-back transactions under a
    single low \code{SS} are not supported, and attempting it is harmless: \textbf{a transaction
    begins only on the first byte after \code{SS} has fallen}, so a sixth byte arriving while
    \code{SS} is still low is discarded rather than taken for a new command byte.
\end{itemize}

\subsection{Setup and hold around \code{SS}}
\label{app:protocol:ss}

Since \code{SS} is synchronised and then edge-detected rather than used directly, it needs a few
system clock cycles to take effect. At 50 MHz one cycle is 20 ns, and the numbers are:

\begin{booktable}{LlL}
  \thead{Parameter} & \thead{Minimum} & \thead{Why} \\
  \midrule
  \code{SS} low before the first rising \code{SCK} edge & \textbf{60 ns} (3 cycles) & the
    frame-start edge is detected on the third system clock edge after \code{SS} falls \\
  \code{SS} high between transactions & \textbf{60 ns} (3 cycles) & the same detector has to see the
    line return before it can be rearmed \\
  \code{SS} low after the last \code{SCK} edge & \textbf{60 ns} (3 cycles) & so that the last byte is
    registered before the frame is torn down \\
\end{booktable}

Allow 100 ns for each of them in a driver and none of this is close. An AVR toggling \code{SS} through
\code{PORTC.OUTCLR} and \code{PORTC.OUTSET} around a byte written to \code{SPI0.DATA} is already far
slower than that, so the limits matter only if the master drives \code{SS} from hardware, or runs much
faster than the AVR32DB28 does.

There is \textbf{no} minimum gap between bytes: bytes may come back to back within a transaction, and
the slave does not care how long the master pauses between them, as long as \code{SS} stays low.

\subsection{\code{MISO} outside a read}
\label{app:protocol:miso}

\begin{itemize}
  \item \code{MISO} is driven \textbf{low} while the slave is deselected, it is not released to high
    impedance. The design assumes a single slave on the bus. Adding a second would require making this
    output tri-stated.
  \item The byte the master clocks out \textbf{during the command byte} is meaningless: the slave
    loads its response shift register at the command byte's \emph{end}, so what leaves during it is
    whatever happened to remain. Discard it.
  \item During a write all four \code{MISO} data bytes are meaningless too.
\end{itemize}

% ------------------------------------------------------------------------------------------
\section{Worked example}
\label{app:protocol:example}

Sending a frame (\code{id 0x123}, \code{dlc 2}, data \code{AA BB}) is five transactions, and then a
poll:

\begin{widedrawing}
Write TX_ID      : 81 00 00 01 23
Write TX_DLC     : 82 00 00 00 02
Write TX_DATA_LO : 83 AA BB 00 00
Write TX_DATA_HI : 84 00 00 00 00
Write TX_SEND    : 85 00 00 00 01
Read  STATUS     : 00 xx xx xx xx   -> MISO gives 00 00 00 00 while the transmission runs,
                                       then 00 00 00 01 once tx_done has set TX ready again.
\end{widedrawing}

Note \code{83 AA BB 00 00}: byte 0 in the most significant position. That is the packing on both sides
of the wire \textemdash{} the driver assembles it, \code{register_bank} presents \code{TX_DATA_LO} as
the \emph{first} four bytes, and \code{can_controller} reads \code{tx_data} with byte 0 in bits
63--56. The diagram in \secref{app:protocol:order} is the picture of that, and
\code{register_bank_tb} checks it.

Worth reading off the same example: the first four \code{MISO} bytes of a write are meaningless, and
so is the \code{MISO} byte returned \emph{during} any command byte. Discard it.

\begin{aside}
\textbf{Changing this contract.} The register map and the protocol are a contract between two classes
that cannot compile against each other. If anything changes, both sides have to know about it before
the change is merged. In practice: changes are discussed with both supervisors first, and the
documents are updated \emph{before} the code, never after.
\end{aside}
